\section{From mode-frame transport to gauge-covariant envelope dynamics}
\label{sec:effective-field}

\subsection{Multiscale ansatz and retained polarization subspace}

Assume an excitation narrowband around a carrier angular frequency $\omega_0$.
Write the substrate displacement as a slowly varying complex envelope multiplied by a local mode frame:
\begin{equation}
  \vecu(x,t)\;\approx\;
  \mathrm{Re}\!\left[
    \e^{-\ii\omega_0 t}\sum_{a=1}^{n}\Psi^a(x,t)\,\mathbf{e}_a(x,t)
  \right],
  \label{eq:multiscale-ansatz}
\end{equation}
where $\Psi(x,t)\in\C^n$ varies slowly compared to $\omega_0^{-1}$ and
$\{\mathbf{e}_a(x,t)\}_{a=1..n}$ is an orthonormal frame spanning the retained polarization subspace.
The frame is not unique: $\mathbf{e}_a\mapsto \sum_b \mathbf{e}_b\,U_{ba}(x)$ with $U(x)\in U(n)$ represents the same
physical subspace.

Projecting the full substrate dynamics onto the retained subspace produces, generically, a gauge-covariant derivative
acting on $\Psi$:
\begin{equation}
  \mathcal{D}_\mu \Psi
  =
  \left(\partial_\mu + \ii\,\mathcal{A}_\mu\right)\Psi,
  \label{eq:covD}
\end{equation}
where $\mathcal{A}_\mu$ is the Berry (rank-1) or Wilczek--Zee (non-Abelian) connection defined in \Cref{sec:geophase}.
The gauge transformation $\mathbf{E}\mapsto \mathbf{E}U$ induces $\Psi\mapsto U^\dagger\Psi$ and the standard
connection transformation law, leaving $\mathcal{D}_\mu\Psi$ covariant.

\subsection{Effective action for the slow sector}

Under the multiscale separation (single-band or near-degenerate subspace) and smoothness of the polarization bundle,
the coarse-grained slow dynamics admits a gauge-invariant effective action built from $\Psi$ and the curvature
$\mathcal{F}_{\mu\nu}$:
\begin{equation}
  S_{\text{eff}}[\Psi,\mathcal{A}]
  \;=\;\int \dd^4x\;
  \Big(
    \Psi^\dagger \mathcal{K}(\mathcal{D}) \Psi
    \;-\;\kappa\,\mathrm{tr}(\mathcal{F}_{\mu\nu}\mathcal{F}^{\mu\nu})
    \;+\;\cdots
  \Big),
  \label{eq:Seff}
\end{equation}
where $\mathcal{K}(\mathcal{D})$ is an operator determined by the linearized carrier band (e.g.\ a Klein--Gordon-type or
Dirac-type envelope operator), $\kappa$ is an emergent stiffness for curvature fluctuations, and ``$\cdots$'' denotes
higher-order and symmetry-allowed terms.
In the Abelian case ($n=1$) the curvature term reduces to $F_{\mu\nu}F^{\mu\nu}$ with $F_{\mu\nu}=f_{\mu\nu}$.
For comparison with electromagnetic conventions one may introduce a dimensionful potential
$A^{\mathrm{eff}}_\mu := \eta\,a_\mu$ (and analogously in the non-Abelian case),
so that $F^{\mathrm{eff}}_{\mu\nu}=\partial_\mu A^{\mathrm{eff}}_\nu-\partial_\nu A^{\mathrm{eff}}_\mu=\eta\,f_{\mu\nu}$.
The conversion factor $\eta$ fixes units and would have to be determined by matching to experiment; it is not introduced
as a fundamental coupling in the microscopic substrate equations.


Geometric phases are universal in wave systems \citep{Pancharatnam1956,Berry1984,WilczekZee1984,Bliokh2015SpinOrbit,Cisowski2022}.
What is specific here is the identification of the physically relevant gauge structure with the
connection induced by a distinguished narrowband sector of the substrate.

\subsection{Dirac-envelope interpretation and the ``split''}

In standard quantum field theory, the Dirac field $\psi$ simultaneously encodes (i) the internal particle labels
(spin, particle/antiparticle) and (ii) the local phase/gauge coupling through the covariant derivative
$D_\mu=\partial_\mu+\ii(q/\hbar)A_\mu$.
In the present theory, the proposal is to reinterpret this mixture as follows:
\begin{itemize}
\item the \emph{particle sector} corresponds to localized solitonic modes with additional internal structure
(topology/geometry and its associated labels);
\item the \emph{phase/gauge sector} corresponds to the Berry/Wilczek--Zee connection induced by transport of the
ordered narrowband polarization subspace.
\end{itemize}
An effective Dirac-type envelope equation is then understood as a compact coupling of these two sectors.
In particular, a ``charge''-like coupling parameter is not introduced as a fundamental constant of a new field;
it appears as an effective conversion factor between the normalization conventions of the envelope field and the
geometric connection extracted from the substrate.
